

\begin{center}
  \vspace*{0.5cm}

  {\Large\textbf{\color{uv-dark}Declaración Responsable de Autoría y Uso de Inteligencia Artificial (IA) en el TFM}}

  \vspace{0.5cm}

  {\LARGE\color{uv-blue}Instituto BME - MIAX}

\end{center}

% ============================================================================
% SECCIÓN 1: HERRAMIENTAS DE IA UTILIZADAS
% ============================================================================
\seccion{Herramientas de IA Utilizadas}

Declaro que en este trabajo he utilizado las siguientes herramientas de IA (si procede):

\begin{enumerate}
  % Se proporcionan ejemplos orientados a programación y modelos cuantitativos
  \item Herramienta y versión: \campo{ChatGPT (GPT-4o)}
  \item Herramienta y versión: \campo{GitHub Copilot}
  \item Herramienta y versión: \campo{Claude 3.5 Sonnet}
  % Añade más herramientas si es el caso
\end{enumerate}

\vspace{0.5cm}

% ============================================================================
% SECCIÓN 2: DIMENSIONES DE USO DE LA IA
% ============================================================================
\seccion{1) Dimensiones de Uso de la IA}

\subtitulo{Instrucciones: Determina el grado de uso de la IA en las siguientes dimensiones. Marca una opción por fila.}

\noindent
\textit{Escala (5 puntos): \textbf{1} = No se ha utilizado IA, \textbf{2} = Uso muy puntual, \textbf{3} = Uso moderado,  \textbf{4} = Uso frecuente, \textbf{5} = Uso intensivo}

\begin{center}
  \small
  \setlength{\tabcolsep}{0.3cm}
  \begin{longtable}{|p{11cm}|c|c|c|c|c|}
    \hline
    \textbf{Dimensión} & \textbf{1} & \textbf{2} & \textbf{3} & \textbf{4} & \textbf{5} \\
    \hline
    \endhead
    \hline
    \endfoot

    % Usa \checkmark para la opción que quieres marcar

    Corrección lingüística: ortografía, gramática, puntuación y coherencia formal del texto &  & \checkmark &  &  & \\
    \hline

    Reescritura y mejora de estilo: reformulación para aumentar claridad, concisión y legibilidad manteniendo el significado &  & \checkmark &  & &  \\
    \hline

    Síntesis de contenido: elaboración de resúmenes a partir de textos extensos o complejos (ej. papers financieros) para identificar ideas clave &  &  &  & \checkmark &  \\
    \hline

    Apoyo en la localización de referencias: propuesta de posibles fuentes (a verificar y citar correctamente) & \checkmark &  &  & & \\
    \hline

    Planificación del trabajo: propuesta de esquemas, índices o guiones iniciales para estructurar el documento &  &  &  & \checkmark &  \\
    \hline

    Mejora de la argumentación y la estructura: detección de huecos, redundancias o saltos lógicos y propuesta de reorganización &  & \checkmark &  &  &  \\
    \hline

    Programación: apoyo en el desarrollo, depuración u optimización de fragmentos de código puntuales (ej. funciones de pandas, numpy) &  &  &  & \checkmark &  \\
    \hline

    Programación: uso generalizado de herramientas de generación de código de forma autónoma a partir de especificaciones &  & \checkmark &  &  &  \\
    \hline

    Análisis de datos: ayuda exploratoria para limpieza preliminar, manejo de NaNs, generación de hipótesis e identificación de patrones &  &  & \checkmark &  &  \\
    \hline

    Traducción o adaptación lingüística (si procede) & \checkmark &  &  &  &  \\
    \hline

    Soporte en material visual: ideas para figuras, tablas o gráficos (ej. curvas de backtesting) y revisión de su explicación &  & \checkmark &  &  &  \\
    \hline

    % Si no se ha usado para otras áreas, la siguiente línea se puede comentar
    Otros (especificar si procede)  & \checkmark &  &  &  &  \\
    \hline

  \end{longtable}
\end{center}


% ============================================================================
% SECCIÓN 3: COMPROMISOS DE USO RESPONSABLE
% ============================================================================
\seccion{2) Compromisos de Uso Responsable}

\subtitulo{Instrucciones: Determina si se ha realizado un uso responsable de la IA. Marca tu grado de acuerdo.}


\noindent
\textit{Escala (acuerdo): \textbf{1} = Totalmente en desacuerdo,  \textbf{2} = En desacuerdo, \textbf{3} = Neutral, \textbf{4} = De acuerdo, \textbf{5} = Totalmente de acuerdo}

\begin{center}
  \small
  \setlength{\tabcolsep}{0.3cm}
  \begin{longtable}{|p{11cm}|c|c|c|c|c|}
    \hline
    \textbf{Declaración} & \textbf{1} & \textbf{2} & \textbf{3} & \textbf{4} & \textbf{5} \\
    \hline
    \endhead
    \hline
    \endfoot

    Las ideas principales, interpretaciones, metodologías cuantitativas y análisis críticos del trabajo son propios. & & & & & \checkmark \\
    \hline

    El uso de IA no ha sustituido mis competencias básicas de programación, análisis financiero y razonamiento académico. &  &  &  & \checkmark &  \\
    \hline

    He revisado, contrastado e integrado críticamente cualquier contenido y código generado con IA antes de incorporarlo. &  &  &  & & \checkmark \\
    \hline

    He verificado la precisión y fiabilidad del contenido, incluyendo evitar el sesgo de anticipación (look-ahead bias) en el código generado. &  &  &  &  & \checkmark \\
    \hline

  \end{longtable}
\end{center}


% ============================================================================
% SECCIÓN 4: OBSERVACIONES
% ============================================================================

% Si no hay observaciones puedes comentar todas las líneas de esta sección
\seccion{Observaciones (Opcional)}

\flushleft

Observación 1: La IA fue utilizada intensivamente para optimizar expresiones regulares y funciones vectorizadas en \texttt{pandas}, evitando bucles \texttt{for} ineficientes en el procesado de series temporales.

\vspace{0.5cm}

Observación 2: Copilot sugirió en repetidas ocasiones implementar modelos predictivos utilizando datos del futuro (*look-ahead bias*). Dichas sugerencias fueron descartadas y reescritas manualmente para mantener la integridad del backtesting.

\newpage


% Estas dos secciones no se deben modificar


% ============================================================================
% SECCIÓN 5: DECLARACIÓN FINAL
% ============================================================================
\seccion{Declaración Final}


\noindent
\fbox{\begin{minipage}{0.98\textwidth}
    \small

  A través de esta declaración, afirmo que he descrito de manera transparente y fidedigna el uso de IA de este trabajo. Asumo la autoría y responsabilidad de la obra final y confirmo que no he introducido datos personales o confidenciales contraviniendo la normativa de protección de datos personales\footnotemark[1] y que he respetado las normas en materia de propiedad intelectual\footnotemark[2], así como las previsiones sobre integridad académica\footnotemark[3] del Instituto BME.

\end{minipage}}

\vspace{0.5cm}
% ============================================================================
% SECCIÓN 6: FIRMA Y FECHA
% ============================================================================
\seccion{Firma y Autorización}

\vspace{1cm}

\noindent
Fecha: 30/06/2026 % Fecha de la convocatoria de ejemplo

\vspace*{1cm}

\noindent
Firma:
% Puedes poner la firma digital aquí. Si no tienes firma digital, entonces imprime esta página,
% firma, escanea e incluye esa página escaneada en lugar de esta en el documento final.

% ============================================================================
% NOTAS A PIE DE PÁGINA
% ============================================================================
\footnotetext[1]{\href{https://www.boe.es/buscar/doc.php?id=BOE-A-2018-16673}{Ley Orgánica 3/2018, de 5 de diciembre, de Protección de Datos Personales y garantía de los derechos digitales}}
\footnotetext[2]{\href{https://www.boe.es/buscar/act.php?id=BOE-A-1996-8930}{Ley de Propiedad Intelectual}}
\footnotetext[3]{Normativa sobre integridad académica del Instituto BME}

\newpage